\section{Implementación}\label{sec:implementacion}

\subsection{Diseño de clases}\label{sec:clases}

\begin{sloppypar}
Las tres clases base conservan todas sus firmas y la única interfaz que creció es la de \texttt{Backtracking}. \texttt{Instancia} guarda $n$, $d$, $k$ y la matriz de características y expone \texttt{feature(i)} y \texttt{costo(i, j)} con índices en base 1 en toda la interfaz pública, así que internamente \texttt{\_features[i-1]} es $f_i$. \texttt{Algoritmo} es la interfaz abstracta con \texttt{resolver(instancia)} y \texttt{nombre()}, que implementan \texttt{FuerzaBruta} y \texttt{Programacion\-Dinamica} sin agregar miembros. \texttt{Solucion} es la lista de pulsos conservados. Un \texttt{git diff} contra el commit inicial confirma que esos tres encabezados no cambiaron.
\end{sloppypar}

La carga de una instancia es atómica y valida todo lo que puede fallar antes de resolver. \texttt{Instancia::cargar} lee por tokens y rechaza con \texttt{std::runtime\_error} los casos $n < 2$, $d < 1$, $k \notin [2, n]$, valores faltantes, no numéricos o no finitos y tokens sobrantes, porque un archivo con más filas que las declaradas resolvería en silencio un prefijo de la grabación. Antes de reservar memoria comprueba que $n \cdot d$ no supere la mitad del tamaño del archivo en bytes, así una cabecera absurda falla al instante en vez de agotar la memoria. Los atributos se sobrescriben recién al final para que una carga fallida no deje una instancia a medias.

El costo se calcula al vuelo con un atajo para el caso sin salto. \texttt{Instancia::costo(i, j)} exige $1 \le i < j \le n$, devuelve $0$ sin tocar los vectores cuando $j = i + 1$ y en el resto de los casos calcula $\lVert f_{i+1} - f_j \rVert_2$ escalando por la mayor diferencia absoluta antes de elevar al cuadrado, como \texttt{std::hypot}. En \texttt{Solucion}, \texttt{costo} devuelve $+\infty$ si algún par consecutivo no cumple $1 \le i < j \le n$, así una selección rota nunca parece mejor que una válida, \texttt{esValida} exige exactamente $k$ índices que empiecen en $1$, terminen en $n$ y crezcan estrictamente y \texttt{guardar} escribe una única línea creando el directorio padre si hace falta.

\texttt{Backtracking} agrega lo necesario para los experimentos sin tocar \texttt{resolver}. El estado de la búsqueda (selección parcial, mejor solución, mejor costo y contador de nodos) vive en atributos privados y la interfaz suma \texttt{usarPodaFactibilidad(bool)}, \texttt{usarPodaOptimalidad(bool)}, sus consultores y \texttt{nodosVisitados()}. Las dos opciones \texttt{-{}-sin-poda-*} del binario se aplican con un \texttt{dynamic\_cast} a \texttt{Backtracking*}, lo que evita agregar métodos virtuales a \texttt{Algoritmo}. Con \texttt{bt} el binario imprime además los nodos visitados y el estado de cada poda.

Cuatro dificultades aparecieron en las revisiones cruzadas del código. La norma sumaba cuadrados sin escalar y con diferencias mayores a $1{,}34 \cdot 10^{154}$ daba $+\infty$, que fuerza bruta y programación dinámica usaban como centinela de ``sin candidato'', así que una devolvía una selección vacía y la otra una que empezaba en $0$. El arreglo combinó la norma escalada con aceptar siempre el primer candidato, de modo que los tres algoritmos devuelven una selección válida aun con costo infinito. \texttt{clang} en \texttt{arm64} con \texttt{-O2} fusiona \texttt{suma += dif*dif} en una instrucción de multiplicación y suma fusionadas (FMA) y la norma difería en una unidad del último dígito (ulp) de la réplica en Python en 17 de 630 pares de prueba, suficiente para invertir una poda o un desempate ante empates exactos, lo que se resolvió compilando con \texttt{-ffp-contract=off}. La recursión de \texttt{extender} tiene profundidad $k$ y con $k$ del orden de $8 \cdot 10^{4}$ agota la pila del hilo principal, un límite documentado en el encabezado y muy por encima del rango del enunciado ($n = 2000$ con $k = 1999$ corre sin problema). Un error de tipeo en \texttt{-{}-sin-poda-*} corría en silencio con ambas podas activas y código de salida $0$ porque el analizador de argumentos ignoraba las opciones desconocidas. Ahora cualquier argumento desconocido u opción sin valor corta con código $1$, igual que una selección inválida, que además no se escribe.

\subsection{Reimplementación en Python y uso de inteligencia artificial}\label{sec:python}

La versión en Python, \texttt{python/algoritmos.py}, replica \texttt{bt} y \texttt{pd} con la misma interfaz de línea de comandos que el binario y sin más que la biblioteca estándar. Lo único que cambia es la etiqueta \texttt{py-bt} o \texttt{py-pd} en la columna \texttt{algoritmo} del CSV, para poder mezclar mediciones de los dos lenguajes. La programación dinámica conserva las dos filas, la tabla desplazada de predecesores y el criterio de desempate del C++. El backtracking conserva la cota inicial, las podas, el orden de exploración y el contador de nodos, pero reemplaza la recursión por una pila explícita de marcos (último, elegidos, costo parcial, próximo $j$), porque el límite de recursión de Python es 1000 por defecto y, aun subiéndolo con \texttt{sys.setrecursionlimit}, con $k$ grande la recursión puede agotar la pila nativa del intérprete sin excepción que capturar.

La equivalencia entre lenguajes es bit a bit y no solo en el costo. Las podas y los desempates comparan \texttt{double}s, así que si la norma en Python redondea distinto que en C++ dos caminos con el mismo costo matemático pueden quedar en orden invertido y la selección y la cantidad de nodos dejan de coincidir, lo que con \texttt{math.dist} ocurría en alrededor del 15\,\% de los pares con $d \ge 2$. Por eso \texttt{costo} en Python replica operación por operación a \texttt{Instancia::costo} (máximo de las diferencias absolutas, suma secuencial de $(\mathit{dif}/\mathit{escala})^2$ y $\mathit{escala} \cdot \sqrt{\mathit{suma}}$) y coincidió con el binario en los 9107 pares de una batería de 240 instancias con $d$ de 1 a 12 y magnitudes de hasta $10^{300}$.

Para la reimplementación se usó Claude, el asistente de Anthropic, a partir del código C++ ya escrito y verificado. Primero se le pidió traducir backtracking y programación dinámica a un único archivo de Python conservando la interfaz de línea de comandos, las podas, la cota inicial y la recurrencia con predecesores. Luego se le pidió la equivalencia exacta con el binario (contador de nodos y norma replicada en lugar de \texttt{math.dist}) y un script de comparación contra el binario, que dejó a la vista la discrepancia por FMA. La versión final se ajustó a mano y se verificó con \texttt{verificar.py}. La misma herramienta se consultó en casos puntuales para depurar el C++ (desborde a infinito de la norma y analizador de argumentos), sin que generara los algoritmos. La conversación está disponible en \url{https://claude.ai/share/fd642169-7146-46ff-8f7d-c7d8d010dc94}.

\subsection{Verificación}\label{sec:testing}

La verificación se apoya en un verificador cruzado que compara los cinco programas contra un óptimo calculado en forma independiente. \texttt{python/verificar.py} genera instancias chicas al azar ($n \le 12$, $d \le 4$, $k \in [2, n]$) en cuatro modos (uniforme, estructura, adversarial y ``repetidos'', con coordenadas en $\{0, 1, 2\}$ para forzar empates exactos) e incluye siempre los bordes $n = 2$, $k = 2$ y $k = n$. Corre cada programa como subproceso, valida la selección y recalcula su costo con una norma escrita aparte. Ese costo se compara con tolerancia $10^{-6}$ contra el mínimo sobre las $\binom{n-2}{k-2}$ selecciones enumeradas en Python y contra el de los demás programas. En 300 casos corridos con los cinco programas a la vez (33 de ellos con $n = 2$ y 78 del modo repetidos) no hubo ninguna discrepancia. Otras 1300 corridas del verificador sobre \texttt{fb}, \texttt{bt} (con cada combinación de podas) y \texttt{pd} tampoco produjeron ninguna. La enumeración exhaustiva limita el verificador a $n \le 12$, de modo que para tamaños mayores la corrección se apoya en la coincidencia de costos entre los tres algoritmos en la experimentación.

La coincidencia entre C++ y Python se verificó también en el árbol de búsqueda. En 30 instancias aleatorias con las cuatro combinaciones de podas (120 corridas) la selección, el costo y la cantidad de nodos fueron iguales en los dos lenguajes. En otras 160 instancias con empates y magnitudes extremas la salida por pantalla y el archivo de selección fueron idénticos en las 800 corridas de \texttt{bt} y \texttt{pd}. Sobre el ejemplo del enunciado los cinco programas devuelven $1\ 2\ 4\ 6$ con costo $2$ y las cuatro cuentas de nodos del backtracking coinciden con las de la sección~\ref{sec:bt} en los dos lenguajes.

\subsection{Compilación y ejecución}\label{sec:compilacion}

El código C++ se compila con \texttt{make} desde \texttt{TP1/}, sin dependencias externas, con \texttt{-std=c++17 -O2 -Wall -Wextra -ffp-contract=off} y sin advertencias. La invocación general es
\begin{verbatim}
./radioedit --instancia <archivo> --algoritmo <fb|bt|pd> [--salida <archivo>]
            [--csv <archivo>] [--sin-poda-factibilidad] [--sin-poda-optimalidad]
\end{verbatim}
\begin{sloppypar}
y escribe la selección en \texttt{-{}-salida} y, si se pide, una fila \texttt{algoritmo,n,d,k,costo,ms} en el CSV. La versión en Python corre con \texttt{python3 python/algoritmos.py} y los mismos argumentos, con Python 3.8 o superior. \texttt{python/generador.py} produce instancias sintéticas, \texttt{python/verificar.py} ejecuta la verificación cruzada y \texttt{python/experimentos.py todo} reproduce toda la experimentación (unos 25 minutos, o menos de un minuto con \texttt{-{}-rapido}). Los otros tres programas usan solo la biblioteca estándar, mientras que \texttt{experimentos.py} necesita \texttt{matplotlib} para las figuras y, para reconstruir audio, los paquetes de \texttt{requirements.txt}, como detalla el \texttt{README.md}.
\end{sloppypar}
